\documentclass[../../main.tex]{subfiles}

\graphicspath{{../../image/}{../../figure/}}

\begin{document}

\section{初边值问题}\label{section:ch05-初边值问题}

双曲守恒律方程组初边值问题的恰当提法与经典解局部存在性的建立，近年来已被深入研究。已出现相当确定性的、尽管高度技术化且复杂的理论，超出本书范围。幸运的是，详细阐述现已见于书籍与综述文章，引于\cref{section:ch05-注释}。
为让读者对该理论的现状有一印象，我们在此按\cref{section:ch04-初边值问题}给出的初边值问题提法路线，记录一个代表性结果。

我们先固定区域为半空间
\begin{align}
\mathscr{D}=\{x\in\mathbb{R}^{m}:\nu\cdot x<0\},\label{align:ch05-halfspace}
\end{align}
外单位法向量 $\nu\in\mathbb{S}^{m-1}$。我们求系统
\begin{align}
\partial_{t}U(x,t)+\sum_{\alpha=1}^{m}\partial_{\alpha}G_{\alpha}(U(x,t))=0,\qquad x\in\mathscr{D},\ t\in(0,T),\label{align:ch05-ibvp-system}
\end{align}
满足初始条件
\begin{align}
U(x,0)=U_{0}(x),\qquad x\in\mathscr{D},\label{align:ch05-ibvp-initial}
\end{align}
以及特殊形式\eqref{align:ch04-边界条件}、\eqref{align:ch04-PBG}的边界条件，即

\begin{align}
BG(U(x,t))\nu=0,\qquad x\in\partial\mathscr{D},\ t\in(0,T),\label{align:ch05-ibvp-boundary}
\end{align}
其中 $B$ 是常值 $n\times n$ 矩阵。

我们对系统\eqref{align:ch05-ibvp-system}作如下假设。通量 $G(U)$ 是定义在 $\overline{\mathscr{B}_{\rho}}$ 上光滑的 $n\times m$ 矩阵值函数。规范化 $G(0)=0$。此外，\eqref{align:ch05-ibvp-system}装备有光滑熵 $\eta(U)$，使 $\mathrm{D}^{2}\eta(U)$ 在 $\overline{\mathscr{B}_{\rho}}$ 上正定。这特别蕴涵\eqref{align:ch05-ibvp-system}双曲，从而对任意 $U\in\overline{\mathscr{B}_{\rho}}$ 与 $\xi\in\mathbb{S}^{m-1}$，由\eqref{align:ch04-方向矩阵}定义的矩阵 $\Lambda(\xi;U)$ 具有实特征值（特征速度）
$\lambda_{1}(\xi;U)\leqslant\cdots\leqslant\lambda_{n}(\xi;U)$ 及相应线性无关特征向量 $R_{1}(\xi;U),\ldots,R_{n}(\xi;U)$。我们要求每个特征值在 $\mathbb{S}^{m-1}\times\overline{\mathscr{B}_{\rho}}$ 上具常重数。

转向边界条件\eqref{align:ch05-ibvp-boundary}，我们引入边界数据"流形"
\begin{align}
\mathscr{M}=\{U\in\overline{\mathscr{B}_{\rho}}:BG(U)\nu=0\}\label{align:ch05-boundary-manifold}
\end{align}
并假设边界\emph{非特征}，即对某 $k=0,\cdots,n$ 及所有 $U\in\mathscr{M}$，
\begin{align}
\lambda_{k}(\nu;U)<0<\lambda_{k+1}(\nu;U),\label{align:ch05-noncharacteristic}
\end{align}
其中 $\lambda_{0}(\nu;U)=-\infty$、$\lambda_{n+1}(\nu;U)=\infty$。于是 $k$ 个特征场经 $\partial\mathscr{D}$ 进入 $\mathscr{D}$，$n-k$ 个特征场离开 $\mathscr{D}$。

我们进一步假设，对任意 $U\in\mathscr{M}$，$B\Lambda(\nu;U)$ 的秩为 $k$，且
\begin{align}
\mathbb{E}^{k}(\nu;U)\oplus\ker[B\Lambda(\nu;U)]=\mathbb{R}^{n},\label{align:ch05-decomposition}
\end{align}
其中 $\mathbb{E}^{k}(\nu;U)$ 表示 $\mathbb{R}^{n}$ 中由 $R_{1}(\nu;U),\ldots,R_{k}(\nu;U)$ 张成的子空间。为说明此条件，我们将系统\eqref{align:ch05-ibvp-system}与边界条件\eqref{align:ch05-ibvp-boundary}关于任一 $U\in\mathscr{M}$ 线性化：
\begin{align}
\partial_{t}V(x,t)+\sum_{\alpha=1}^{m}\mathrm{D}G_{\alpha}(U)\partial_{\alpha}V(x,t)=0,\qquad x\in\mathscr{D},\ t\in(0,T),\label{align:ch05-ibvp-linear}
\end{align}
\begin{align}
B\Lambda(\nu;U)V(x,t)=0,\qquad x\in\partial\mathscr{D},\ t\in(0,T).\label{align:ch05-ibvp-linear-boundary}
\end{align}
粗略地说，\eqref{align:ch05-decomposition}的作用是保证 $V$ 在 $\partial\mathscr{D}$ 上的迹由边界条件结合 $n-k$ 个出射特征场带到边界的信息联合确定。

关于边界条件的最后假设是\emph{一致 Kreiss--Lopatinski 条件}，表述如下。对每个 $U\in\mathscr{M}$、切于边界的向量 $\xi\in\mathbb{S}^{m-1}$（即 $\xi\cdot\nu=0$）、以及具 $\operatorname{Re}z>0$ 的复数 $z$，我们定义矩阵
\begin{align}
M(z,\xi;U)=\Lambda(\nu;U)^{-1}[zI+i\Lambda(\xi;U)].\label{align:ch05-M-def}
\end{align}
我们以 $\mathbb{E}(z,\xi;U)$ 表示 $\mathbb{R}^{n}$ 中由 $M(z,\xi;U)$ 具负实部特征值对应的特征向量张成的子空间，并要求
\begin{align}
|W|\leqslant c|B\Lambda(\nu;U)W|,\qquad\text{对所有 }W\in\mathbb{E}(z,\xi;U),\label{align:ch05-Lopatinski}
\end{align}
其中 $c$ 是不依赖于 $U,\xi,z$ 的正常数。为解释此假设，注意线性系统\eqref{align:ch05-ibvp-linear}有形如
\begin{align}
V(x,t)=\exp(i\xi\cdot x+zt)W(\nu\cdot x)\label{align:ch05-boundary-layer}
\end{align}
的解，其中函数 $W(\tau)$ 满足常微分方程
\begin{align}
\dot{W}+M(z,\xi;U)W=0.\label{align:ch05-W-ode}
\end{align}
\eqref{align:ch05-Lopatinski}的作用是排除满足边界条件\eqref{align:ch05-ibvp-linear-boundary}且在空间方向呈"温和"增长但随时间指数增长的解\eqref{align:ch05-boundary-layer}。

最后转向初始条件\eqref{align:ch05-ibvp-initial}。对 $j=0,1,\ldots$，令 $H_{j}$ 表示 Sobolev 空间 $W^{j,2}(\mathscr{D};\mathbb{R}^{n})$，并设 $U_{0}\in H_{\ell}$，某 $\ell>\frac{m}{2}+1$。然后可由\eqref{align:ch05-ibvp-system}形式地算出解的时间导数初值 $U_{1}(x),\cdots,U_{\ell-1}(x)$，即 $\partial_{t}U(x,0),\cdots,\partial_{t}^{\ell-1}U(x,0)$。于是
\begin{align}
U_{1}=-\sum_{\alpha=1}^{m}\mathrm{D}G_{\alpha}(U_{0})\partial_{\alpha}U_{0},\label{align:ch05-U1}
\end{align}
\begin{align}
U_{2}=-\sum_{\alpha=1}^{m}\mathrm{D}G_{\alpha}(U_{0})\partial_{\alpha}U_{1}-\sum_{\alpha=1}^{m}\mathrm{D}^{2}G_{\alpha}(U_{0})[U_{1},\partial_{\alpha}U_{0}],\label{align:ch05-U2}
\end{align}
等等。此外，$U_{j}\in H_{\ell-j}$，$j=0,\cdots,\ell-1$。特别地，$U_{j}$ 在超平面 $\partial\mathscr{D}$ 上的迹对 $j=0,\cdots,\ell-1$ 良定义。我们进而要求初值与边界条件相容，即
\begin{align}
B\partial_{t}^{j}G(U(x,t))\nu=0,\qquad t=0,\ x\in\partial\mathscr{D},\ j=0,\cdots,\ell-1,\label{align:ch05-compatibility}
\end{align}
亦即
\begin{align}
BG(U_{0}(x))\nu=0,\qquad x\in\partial\mathscr{D},\label{align:ch05-compat0}
\end{align}
\begin{align}
B\Lambda(\nu;U_{0}(x))U_{1}(x)=0,\qquad x\in\partial\mathscr{D},\label{align:ch05-compat1}
\end{align}
等等。

我们现已为陈述存在性定理作好了准备：

\begin{theorem}\label{theorem:ch05-ibvp-existence}
在上述关于系统、边界条件与初值的假设下，存在初边值问题\eqref{align:ch05-ibvp-system}、\eqref{align:ch05-ibvp-initial}、\eqref{align:ch05-ibvp-boundary}的唯一经典解 $U\in C^{1}(\mathscr{D}\times[0,T_{\infty}))$，对某 $0<T_{\infty}\leqslant\infty$。此外，
\begin{align}
U(\cdot,t)\in\bigcap_{j=0}^{\ell}C^{j}([0,T_{\infty});H_{\ell-j}).\label{align:ch05-ibvp-regularity}
\end{align}
区间 $[0,T_{\infty})$ 是极大的，即若 $T_{\infty}<\infty$，则
\begin{align}
\limsup_{t\to T_{\infty}}\|\|\nabla U(\cdot,t)\|\|_{L^{\infty}}=\infty\label{align:ch05-ibvp-blowup1}
\end{align}
且/或
\begin{align}
\limsup_{t\to T_{\infty}}\|U(\cdot,t)\|_{L^{\infty}}=\rho.\label{align:ch05-ibvp-blowup2}
\end{align}
\end{theorem}

（冗长而技术化的）证明从常系数线性系统出发，过渡到变系数线性系统，再经由线性化\eqref{align:ch05-ibvp-linear}、\eqref{align:ch05-ibvp-linear-boundary}与对映射 $U\mapsto V$ 的不动点论证过渡到拟线性系统。

应注意，上述定理的假设对处理连续介质物理中出现的许多自然初边值问题过强。在 Euler 方程（等熵或非等熵气体流动）中，特征速度具常重数的假设确实成立（见\cref{section:ch03-例子}）；
但边界非特征的假设常被违反，例如在无穿透（或滑移）边界条件 $v\cdot\nu=0$ 的情形。在等熵或非等熵弹性动力学方程中，特征速度具常重数的条件常被违反，例如在各向同性弹性固体自然态附近，剪切波对应的特征速度重数发生过渡。此外，边界总是特征的，因为系统具零特征速度。除此之外，还需考虑更一般的区域 $\mathscr{D}$ 以及 $\partial\mathscr{D}$ 上比\eqref{align:ch05-ibvp-boundary}更一般形式的齐次或非齐次边界条件。这些问题由\cref{theorem:ch05-ibvp-existence}的更精致版本处理。文献引于\cref{section:ch05-注释}。

与其诉诸上述一般理论，不如对连续介质物理中出现的特定守恒律方程组初边值问题从头处理，利用其特殊结构，常更有利。作为说明性例子，我们考虑\eqref{align:ch03-等熵热弹性系统}，
它描述弹性固体在 Lagrangian 坐标下的等熵运动，为简单起见假设体积力为零，于是
\begin{align}
\begin{cases}
\partial_{t}F_{i\alpha}-\partial_{\alpha}v_{i}=0,\\
\partial_{t}v_{i}-\partial_{\alpha}S_{i\alpha}(F)=0,
\end{cases}\qquad i,\alpha=1,2,3.\label{align:ch05-elastic-ibvp}
\end{align}
在\eqref{align:ch05-elastic-ibvp}及后续中，我们采用求和约定。

如\cref{section:ch03-例子}所述，Piola--Kirchhoff 应力 $S$ 由势\eqref{align:ch03-等熵热弹性本构}导出，

从而 $\eta=\varepsilon(F)+\frac{1}{2}|v|^{2}$ 是系统的熵。我们假设内能函数 $\varepsilon(F)$ 满足 $\varepsilon(F)\geqslant0$、$\varepsilon(I)=0$ 且秩一凸，即
\begin{align}
A_{i\alpha j\beta}(F)v_{\alpha}v_{\beta}\xi_{i}\xi_{j}>0,\qquad\text{对所有 }\nu,\xi\in\mathbb{S}^{2},\label{align:ch05-rank-one-convex}
\end{align}
其中
\begin{align}
A_{i\alpha j\beta}(F)=\frac{\partial S_{i\alpha}(F)}{\partial F_{j\beta}}=\frac{\partial^{2}\varepsilon(F)}{\partial F_{i\alpha}\partial F_{j\beta}}.\label{align:ch05-A-elastic}
\end{align}
注意重要对称关系
\begin{align}
A_{i\alpha j\beta}(F)=A_{j\beta i\alpha}(F).\label{align:ch05-A-symmetry}
\end{align}
我们还定义
\begin{align}
B_{i\alpha j\beta k\gamma}(F)=\frac{\partial A_{i\alpha j\beta}(F)}{\partial F_{k\gamma}}=\frac{\partial^{3}\varepsilon(F)}{\partial F_{i\alpha}\partial F_{j\beta}\partial F_{k\gamma}}.\label{align:ch05-B-elastic}
\end{align}

弹性体的参考位形是 $\mathbb{R}^{3}$ 的有界区域 $\Omega$，具光滑边界 $\Gamma$，假设被夹持，从而边界条件为
\begin{align}
v(x,t)=0,\qquad x\in\Gamma,\ t>0.\label{align:ch05-clamped}
\end{align}
我们还赋初值
\begin{align}
F(x,0)=F^{0}(x),\qquad v(x,0)=v^{0}(x),\qquad x\in\Omega,\label{align:ch05-elastic-initial}
\end{align}
假设
\begin{align}
F^{0}\in W^{m-1,2}(\Omega;\mathbb{M}^{3\times3}),\qquad v^{0}\in W^{m,2}(\Omega;\mathbb{R}^{3}),\label{align:ch05-regularity-data}
\end{align}
$m\geqslant4$。对 $r=0,1,\cdots,m$，我们以 $\|\cdot\|_{r}$ 与 $\|\|\cdot\|\|_{r}$ 分别记 Sobolev 空间 $W^{r,2}(\Omega;\mathbb{R}^{3})$ 与 $W^{r,2}(\Omega;\mathbb{M}^{3\times3})$ 的范数。$\overset{(r)}{v}$ 记 $v$ 的 $r$ 阶时间导数 $\partial_{t}^{r}v$。由\eqref{align:ch05-regularity-data}，可将\eqref{align:ch05-elastic-ibvp}与\eqref{align:ch05-elastic-initial}结合，导出 $v$ 时间导数值的候选 $\overset{(r)}{v}^{r}\in W^{m-r,2}(\Omega;\mathbb{R}^{3})$：
\begin{align}
v_{i}^{1}(x)=\partial_{\alpha}S_{i\alpha}(F^{0}(x)),\qquad i=1,2,3,\ x\in\Omega,\label{align:ch05-v1}
\end{align}
\begin{align}
v_{i}^{2}(x)=\partial_{\alpha}[A_{i\alpha j\beta}(F^{0}(x))\partial_{\beta}v_{j}^{0}(x)],\qquad i=1,2,3,\ x\in\Omega,\label{align:ch05-v2}
\end{align}
等等。在此设定下：

\begin{theorem}\label{theorem:ch05-clamped-elastic}
假设初值 $(F^{0},v^{0})$ 与边界条件\eqref{align:ch05-clamped}相容，即 $v^{r}\in W_{0}^{1,2}(\Omega;\mathbb{R}^{3})$，$r=0,\cdots,m-1$。则在时间区间 $[0,T]$ 上存在\eqref{align:ch05-elastic-ibvp}、\eqref{align:ch05-clamped}、\eqref{align:ch05-elastic-initial}的唯一经典解 $(F,v)$，且
\begin{align}
v(\cdot,t)\in\bigcap_{r=0}^{m}C^{r}([0,T];W^{m-r,2}(\Omega;\mathbb{R}^{3})).\label{align:ch05-clamped-regularity}
\end{align}
\end{theorem}

$F$ 及其时间导数的正则性可由\eqref{align:ch05-elastic-ibvp}第一个方程与\eqref{align:ch05-clamped-regularity}联合推出。解的物理可容许性所必需的条件 $\det F>0$，可由取 $\det F^{0}>0$ 并使 $T$ 充分小来保证。

上述命题的详细证明见\cref{section:ch05-注释}引用的文献。
为当前目的，正式、纲要性地推导关键估计即可。我们的策略是借助高阶导数的 $L^{2}$ 界来监控解及其导数的逐点行为。于是假设 $(F,v)$ 是\eqref{align:ch05-elastic-ibvp}、\eqref{align:ch05-clamped}、\eqref{align:ch05-elastic-initial}在时间区间 $[0,T]$ 上的解，具\cref{theorem:ch05-clamped-elastic}中所述性质。目标是证明当 $T$ 充分小时，泛函
\begin{align}
E(t)=\sum_{r=0}^{m}\|\overset{(r)}{v}(\cdot,t)\|_{m-r}\label{align:ch05-E-def}
\end{align}
在 $[0,T]$ 上有界。置
\begin{align}
M=\|\|F^{0}(\cdot)\|\|_{m-1},\qquad N=\sum_{r=0}^{m-1}\|v^{r}(\cdot)\|_{m-r-1},\label{align:ch05-MN}
\end{align}
我们在如下假设下工作：
\begin{align}
\|\|F(\cdot,t)\|\|_{m-1}\leqslant2M,\qquad\sum_{r=0}^{m-1}\|\overset{(r)}{v}(\cdot,t)\|_{m-r-1}\leqslant2N,\qquad t\in[0,T],\label{align:ch05-ansatz}
\end{align}
它将事后验证。

前几节中通过对空间变量微分系统方程来建立此类估计的途径，此处不能使用，因为空间导数的边界行为事先未知。我们仅知道，由\eqref{align:ch05-clamped}，任意阶的 $v$ 的时间导数在 $\Gamma$ 上为零。我们将在估计较简单的泛函
\begin{align}
G(t)=\sum_{r=0}^{m-1}\|\overset{(r)}{v}(\cdot,t)\|_{1}+\|\overset{(m)}{v}(\cdot,t)\|_{0}\label{align:ch05-G-def}
\end{align}
时利用这一点。

对 $r=0,\cdots,m-1$，对\eqref{align:ch05-elastic-ibvp}作用 $\partial_{t}^{r+1}$，得
\begin{align}
\overset{(r+2)}{v}_{i}(x,t)=\partial_{\alpha}[A_{i\alpha j\beta}(F(x,t))\partial_{\beta}\overset{(r)}{v}_{j}(x,t)]+Z_{i}^{r}(x,t),\label{align:ch05-vr2}
\end{align}
其中
\begin{align}
Z_{i}^{r}(x,t)=\partial_{\alpha}[\partial_{t}^{r+1}S_{i\alpha}(F(x,t))-A_{i\alpha j\beta}(F(x,t))\partial_{\beta}\overset{(r)}{v}_{j}(x,t)].\label{align:ch05-Zr-def}
\end{align}

我们以 $2\overset{(r+1)}{v}_{i}$ 乘\eqref{align:ch05-vr2}，在 $\Omega\times[0,t]$ 上积分。回忆对称条件\eqref{align:ch05-A-symmetry}，分部积分一次后，
\begin{align}
\int_{\Omega}[|\overset{(r+1)}{v}(x,t)|^{2}+A_{i\alpha j\beta}(F(x,t))\partial_{\alpha}\overset{(r)}{v}_{i}(x,t)\partial_{\beta}\overset{(r)}{v}_{j}(x,t)]\,\mathrm{d}x
&=\int_{\Omega}[|v^{r+1}(x)|^{2}+A_{i\alpha j\beta}(F^{0}(x))\partial_{\alpha}v_{i}^{r}(x)\partial_{\beta}v_{j}^{r}(x)]\,\mathrm{d}x\nonumber\\
&\quad+\int_{0}^{t}\int_{\Omega}[2Z_{i}^{r}(x,\tau)\overset{(r+1)}{v}_{i}(x,\tau)+B_{i\alpha j\beta k\gamma}(F)\partial_{\alpha}\overset{(r)}{v}_{i}(x,\tau)\partial_{\beta}\overset{(r)}{v}_{j}(x,\tau)\partial_{\gamma}v_{k}(x,\tau)]\,\mathrm{d}x\,\mathrm{d}\tau.\label{align:ch05-energy-elastic}
\end{align}
我们以总能量（熵）守恒律补充\eqref{align:ch05-energy-elastic}：
\begin{align}
\int_{\Omega}[\varepsilon(F(x,t))+\frac{1}{2}|v(x,t)|^{2}]\,\mathrm{d}x=\int_{\Omega}[\varepsilon(F^{0}(x))+\frac{1}{2}|v^{0}(x)|^{2}]\,\mathrm{d}x.\label{align:ch05-total-energy}
\end{align}

为避免符号泛滥，我们约定：下文 $a,b,K$ 表示一般正常数，$a$ 依赖于 $M$，$b$ 依赖于 $M$ 与 $N$，$K$ 依赖于 $\|\|F^{0}\|\|_{m-1}$ 与 $\|v^{0}\|_{m}$。特别地，对 $r=0,\cdots,m$，$\|v^{r}\|_{m-r}\leqslant K$。

由\eqref{align:ch05-rank-one-convex}且 $\overset{(r)}{v}(\cdot,t)\in W_{0}^{1,2}(\Omega;\mathbb{R}^{3})$，Gårding 不等式给出：
\begin{align}
\int_{\Omega}A_{i\alpha j\beta}(F(x,t))\partial_{\alpha}\overset{(r)}{v}_{i}(x,t)\partial_{\beta}\overset{(r)}{v}_{j}(x,t)\,\mathrm{d}x\geqslant\mu\|\overset{(r)}{v}(\cdot,t)\|_{1}^{2}-a\|\overset{(r)}{v}(\cdot,t)\|_{0}^{2},\label{align:ch05-garding}
\end{align}
$\mu>0$。此外，因 $m\geqslant4$，Sobolev 空间的标准微积分不等式给出
\begin{align}
\|Z^{r}(\cdot,t)\|_{0}\leqslant bE(t),\qquad r=1,\cdots,m-1.\label{align:ch05-Zr-bound}
\end{align}

合并\eqref{align:ch05-energy-elastic}与\eqref{align:ch05-total-energy}、\eqref{align:ch05-garding}、\eqref{align:ch05-Zr-bound}，得形如
\begin{align}
G(t)\leqslant K+b\int_{0}^{t}E(\tau)\,\mathrm{d}\tau.\label{align:ch05-G-bound}
\end{align}

为估计 $\overset{(r)}{v}$ 一阶以上的空间导数，我们将\eqref{align:ch05-vr2}视为强椭圆系统
\begin{align}
A_{i\alpha j\beta}(x,t)\partial_{\alpha}\partial_{\beta}\overset{(r)}{v}_{j}(x,t)=\overset{(r+2)}{v}_{i}(x,t)-Y_{i}^{r}(x,t),\label{align:ch05-elliptic}
\end{align}
其中
\begin{align}
Y_{i}^{r}(x,t)=Z_{i}^{r}(x,t)+B_{i\alpha j\beta k\gamma}(F(x,t))\partial_{\alpha}F_{k\gamma}(x,t)\partial_{\beta}\overset{(r)}{v}_{j}(x,t).\label{align:ch05-Yr-def}
\end{align}
对 $s=0,\cdots,m-r-2$，我们采用椭圆系统理论的标准估计：
\begin{align}
\|\overset{(r)}{v}(\cdot,t)\|_{s+2}\leqslant a\left[\|\overset{(r+2)}{v}(\cdot,t)\|_{s}+\|\overset{(r)}{v}(\cdot,t)\|_{s}+\|Y^{r}(\cdot,t)\|_{s}\right].\label{align:ch05-elliptic-est}
\end{align}
由 Sobolev 空间的微积分不等式，
\begin{align}
\|Y^{r}(\cdot,t)\|_{s}\leqslant bE(t),\qquad\|\partial_{t}Y^{r}(\cdot,t)\|_{s}\leqslant bE(t),\qquad r=1,\cdots,m-2.\label{align:ch05-Yr-bound}
\end{align}
因此
\begin{align}
\|Y^{r}(\cdot,t)\|_{s}\leqslant K+b\int_{0}^{t}E(\tau)\,\mathrm{d}\tau.\label{align:ch05-Yr-Gronwall}
\end{align}

对 $r=0,\cdots,m-2$，借助\eqref{align:ch05-elliptic-est}可递归地用 $\|\overset{(r)}{v}\|_{0},\cdots,\|\overset{(m)}{v}\|_{0}$（当 $m-r$ 偶）或 $\|\overset{(r)}{v}\|_{1},\cdots,\|\overset{(m-1)}{v}\|_{1}$（当 $m-r$ 奇）估计 $\|\overset{(r)}{v}\|_{m-r}$，这些都是 $G$ 的组成部分。于是我们得到形如
\begin{align}
E(t)\leqslant b\left[K+G(t)+\int_{0}^{t}E(\tau)\,\mathrm{d}\tau\right].\label{align:ch05-E-est}
\end{align}

合并\eqref{align:ch05-E-est}与\eqref{align:ch05-G-bound}，得 $E(t)\leqslant K\exp(bt)$，这证明只要假设\eqref{align:ch05-ansatz}成立，$E(t)$ 在 $[0,T]$ 上有界。而这确为事实，因为
\begin{align}
F_{i\alpha}(x,t)=F_{i\alpha}^{0}(x)+\int_{0}^{t}\partial_{\alpha}v_{i}(x,\tau)\,\mathrm{d}\tau,\qquad\overset{(r)}{v}_{i}(x,t)=v_{i}^{r}(x)+\int_{0}^{t}\overset{(r+1)}{v}_{i}(x,\tau)\,\mathrm{d}\tau,\label{align:ch05-integral-rel}
\end{align}
给出
\begin{align}
\|\|F(\cdot,t)\|\|_{m-1}\leqslant M+K(e^{bt}-1),\qquad\sum_{r=0}^{m-1}\|\overset{(r)}{v}(\cdot,t)\|_{m-r-1}\leqslant N+K(e^{bt}-1),\label{align:ch05-a-posteriori}
\end{align}
从而当 $T$ 充分小时反推出\eqref{align:ch05-ansatz}。

为完整严格证明\cref{theorem:ch05-clamped-elastic}，可在适当函数类中构造映射 $\Phi\mapsto(F,v)$ 的不动点，其中 $(F,v)$ 是线性系统
\begin{align}
\begin{cases}
\partial_{t}F_{i\alpha}-\partial_{\alpha}v_{i}=0,\\
\partial_{t}v_{i}-A_{i\alpha j\beta}(\Phi)\partial_{\alpha}F_{j\beta}=0,
\end{cases}\qquad i,\alpha=1,2,3,\label{align:ch05-linear-system}
\end{align}
在边界条件\eqref{align:ch05-clamped}与初始条件\eqref{align:ch05-elastic-initial}下的解。或者，可采用粘性消去法，在同样的边界与初始条件下，将\eqref{align:ch05-elastic-ibvp}、\eqref{align:ch05-clamped}、\eqref{align:ch05-elastic-initial}的解作为系统
\begin{align}
\begin{cases}
\partial_{t}F_{i\alpha}-\partial_{\alpha}v_{i}=0,\\
\partial_{t}v_{i}-\partial_{\alpha}S_{i\alpha}(F)=\varepsilon\Delta v_{i},
\end{cases}\qquad i,\alpha=1,2,3,\label{align:ch05-vanishing-viscosity}
\end{align}
的解当 $\varepsilon\to0$ 时的极限。无论哪种途径，关键任务都是证明对\eqref{align:ch05-linear-system}在 $\Phi$ 一致下、或对\eqref{align:ch05-vanishing-viscosity}在 $\varepsilon$ 一致下，泛函 $E(t)$ 有界。

除非被夹持，运动物体的边界随时间变化。因此，在 Euler 坐标中跟踪连续介质的演化常导致具自由边界的初边值问题，自由边界须作为解的一部分确定。这通常引起严重的技术困难。可通过切换到 Lagrangian 坐标来规避该障碍，此时区域变为物体固定的参考位形。然而代价是，Lagrangian 形式的方程通常复杂且缺乏 Euler 对应物的对称性。特别地，气体动力学 Euler 方程即如此。

作为说明，我们考虑被气团等熵膨胀，气团在 $t=0$ 时被约束在 $\mathbb{R}^{3}$ 的有界区域中，周围为真空。该设定中气体的演化已在\cref{section:ch04-消失黏性方法}中建模为 Euler 方程\eqref{align:ch04-Euler}的 Cauchy 问题，

将真空可视作零密度气体。已指出，只要焦点在经典解上，宜将状态变量从密度 $\rho$ 换为加权声速 $\omega$，因为与 Euler 方程\eqref{align:ch04-Euler非守恒}不同，对称系统\eqref{align:ch04-Euler对称双曲}即使在真空态也双曲。

当初值 $(\omega_{0},v_{0})$ 延拓到整个 $\mathbb{R}^{3}$ 且落在 Sobolev 空间 $W^{3,2}$ 中时，\cref{theorem:ch05-凸熵局部存在}保证\eqref{align:ch04-Euler对称双曲}、从而 Euler 方程\eqref{align:ch04-Euler非守恒}的 Cauchy 问题经典解至少局部时间存在。

然而，如\cref{section:ch04-Euler方程}已讨论，该经典解的寿命一般有限，
因为气团内部波破裂或气-真空界面上发展出奇性。波破裂关联着从经典解到含激波弱解的过渡，超出本讨论范围；本讨论将聚焦于界面奇性。

界面奇性的性质已在\cref{section:ch04-Euler方程}讨论。
平衡律迫使压力从而密度在界面上连续。已论证，编码于\eqref{align:ch04-界面加速度}的自然条件是声速平方的法向导数须越过界面跳跃。
为检验解即使在界面奇性出现后是否仍可在内部保持光滑，人们经由初值在 $t=0$ 引入此类奇性，并试图为所得初边值问题构造经典解。在此设定下，界面须视为自由边界，作为解的一部分确定。该方向的一个结果如下。

\begin{theorem}\label{theorem:ch05-vacuum}
考虑理想气体的等熵流动，状态方程 $p=\kappa\rho^{\gamma}$，$1<\gamma\leqslant2$。气体在 $t=0$ 占据 $\mathbb{R}^{3}$ 中具光滑边界 $\Gamma$ 的有界集 $\Omega$，周围为真空。$t=0$ 时气体密度为 $\rho_{0}\in C^{\infty}(\overline{\Omega})$，满足
\begin{align}
\rho_{0}(x)>0,\qquad x\in\Omega,\qquad\rho_{0}(x)=0,\qquad x\in\Gamma,\label{align:ch05-density-data}
\end{align}
\begin{align}
\rho_{0}^{\gamma-1}(x)\geqslant A\operatorname{dist}(x,\Gamma),\qquad x\in\Omega,\label{align:ch05-density-degeneracy}
\end{align}
$A>0$。初速度为 $v_{0}\in C^{\infty}(\Omega;\mathbb{R}^{3})$。则存在 $T>0$ 及定义在 $\overline{\Omega}\times[0,T]$、取值于 $\mathbb{R}^{3}$ 的 $C^{2}$ 函数 $\chi$，具有如下性质：(a) $\chi(\cdot,0)$ 是 $\Omega$ 上的恒等映射；(b) 对任意 $t\in[0,T]$，$\chi(\cdot,t)$ 微分同胚地将 $\Omega$ 映为开集 $\Omega_{t}$，并将 $\Gamma$ 映为 $\Omega_{t}$ 的边界 $\Gamma_{t}$；(c) $\chi$ 诱导 Euler 方程\eqref{align:ch04-Euler非守恒}在 $\{(x,t):t\in[0,T],x\in\overline{\Omega}_{t}\}$ 上唯一 $C^{1}$ 解 $(\rho,v)$，
满足初始条件
\begin{align}
\rho(x,0)=\rho_{0}(x),\qquad v(x,0)=v_{0}(x),\qquad x\in\Omega,\label{align:ch05-vacuum-initial}
\end{align}
及边界条件
\begin{align}
p(\rho(x,t))=0,\qquad v(x,t)\cdot v(x,t)=V(x,t),\qquad t\in[0,T],\ x\in\Gamma_{t},\label{align:ch05-vacuum-boundary}
\end{align}
其中 $\nu$ 表示 $\Gamma_{t}$ 的外单位法向量，$V$ 表示曲面 $\Gamma_{t}$ 的法向速度。
\end{theorem}

界面上奇性的性质编码在\eqref{align:ch05-density-degeneracy}中——回忆声速平方正比于 $\rho^{\gamma-1}$。对 $\Omega$ 中初值施加的正则性限制过强。在弱得多的假设（初值在充分高阶 Sobolev 空间中并满足相当笨拙的相容条件）下解存在。

上述命题的证明见\cref{section:ch05-注释}引用的文献，冗长而技术化。
下面的注记可让读者对方法论有一印象。

为克服自由边界带来的障碍，我们在 Lagrangian 坐标中重新表述问题。我们采用\cref{section:ch02-运动学}引入的记号约定，
略加修改以与 Euler 方程的\cref{section:ch04-Euler方程}记号兼容。
于是物质质点由 $y$（Lagrangian 坐标）标识，而物理空间中的位置由 $x$（Euler 坐标）标识。对任一特定场，我们用同一符号表示其 Lagrangian 表示（作为 $(y,t)$ 的函数）与 Euler 表示（作为 $(x,t)$ 的函数）。为避免混淆，物质时间导数（保持 $y$ 固定）用上圆点号表示，而空间时间导数（保持 $x$ 固定）用 $\partial_{t}$ 表示。

我们选物理空间 $\mathbb{R}^{3}$ 为参考空间，气体在 $t=0$ 时的物理占据 $\Omega$ 为参考位形。于是参考密度为 $\rho_{0}(y)$。定理陈述中出现的函数 $\chi(y,t)$ 表示气体的运动。于是 $\chi(y,0)=y$，对所有 $y\in\Omega$。运动诱导变形梯度场 $F(y,t)$，由 $F=\nabla\chi$，及 Lagrangian 坐标下的速度场 $v(y,t)=\dot{\chi}$。由\eqref{align:ch02-2-3-3}，Lagrangian 坐标下的密度场为 $\rho(y,t)=\rho_{0}(y)J^{-1}(y,t)$，其中 $J=\det F$。

在 Lagrangian 坐标中，Euler 方程呈\eqref{align:ch02-2-3-4}形式，即当体积力为零时，
\begin{align}
\rho_{0}(y)\dot{v}(y,t)=\operatorname{Div}S(F(y,t)),\label{align:ch05-lagrangian}
\end{align}
其中 $S$ 是 Piola--Kirchhoff 应力。对理想气体，Cauchy 应力张量为 $-\kappa\rho^{\gamma}I$，从而由\eqref{align:ch02-2-3-6}，
\begin{align}
S=-\kappa\rho^{\gamma}J(F^{-1})^{\top}=-\kappa\rho(F^{*})^{\top},\label{align:ch05-pk-stress}
\end{align}
其中 $F^{*}$ 是 $F$ 的伴随矩阵，从而 $(F^{*})^{\top}$ 是 $F$ 的余子矩阵。由\eqref{align:ch02-2-2-13}，
$F$ 的余子矩阵的散度为零，由此
\begin{align}
\operatorname{Div}S(F(y,t))=-\kappa(F^{*})^{\top}(y,t)\operatorname{Grad}[\rho_{0}^{\gamma}(y)J^{-\gamma}(y,t)].\label{align:ch05-div-s}
\end{align}

\eqref{align:ch05-lagrangian}、\eqref{align:ch05-div-s}的解由粘性消去法获得，即系统
\begin{align}
\rho_{0}(y)\dot{v}(y,t)=\operatorname{Div}S(F(y,t))+\varepsilon\operatorname{Div}\dot{S}(F(y,t))\label{align:ch05-viscoelastic}
\end{align}
的解当 $\varepsilon\to0$ 时的极限。在\eqref{align:ch05-viscoelastic}中可识别出粘弹性材料的运动方程，其 Piola--Kirchhoff 应力带\cref{section:ch02-热粘弹性}引入的特殊本构方程。

广义地说，此处解的构造遵循早先处理夹持边界弹性固体（\cref{theorem:ch05-clamped-elastic}）的路线，通过建立解导数的 $L^{2}$ 估计（提取自"能量"积分）进行。这些结合量
\begin{align}
E(t)=\sum_{k=0}^{4}\|\overset{(2k)}{\chi}(t)\|_{4-k}^{2},\label{align:ch05-E-vacuum}
\end{align}
以及涡量场某些导数的 $L^{2}$ 估计。边界处的退化使推导这些估计的任务比夹持边界情形困难得多。特别地，某些估计与系统的非线性结构紧密耦合。因此不能类比\eqref{align:ch05-linear-system}经由线性化构造解。粘性消去法的成功取决于粘性项 $\dot{S}$ 的特定选择，它保留 Euler 方程 Piola--Kirchhoff 应力的精细特征。估计的冗长艰苦推导细节记录在\cref{section:ch05-注释}引用的文献中，尽管仅对 $\Omega=\mathbb{T}^{2}\times(0,1)$ 的情形。

\end{document}
